\documentclass[10pt,conference]{IEEEtran}

% --- Packages ---
\usepackage{cite}
\usepackage{graphicx}
\graphicspath{{figs/}} % Automatically look for images in the figs/ folder
\usepackage[cmex10]{amsmath}
\interdisplaylinepenalty=2500
\usepackage{amsthm}
\newtheorem{definition}{Definition}
\usepackage{algorithmic}
\usepackage{array}
\usepackage[caption=false,font=footnotesize]{subfig}
\usepackage{url}

% --- Custom Hyphenation ---
\hyphenation{op-tical net-works semi-conduc-tor}

\begin{document}

% --- Title & Authors ---
\title{Example of IEEEtran.cls, adapted for Sibgrapi 2020}

\newif\iffinal
\finalfalse
%\finaltrue
\newcommand{\cmtid}{99999}

\iffinal
  \author{
    \IEEEauthorblockN{First Author}
    \IEEEauthorblockA{Department of Computer Science\\
    University Name\\
    City, Country\\
    Email: author1@example.com}
    \and
    \IEEEauthorblockN{Second Author}
    \IEEEauthorblockA{Department of Computer Science\\
    University Name\\
    City, Country\\
    Email: author2@example.com}
  }
\else
  \author{Sibgrapi paper ID: \cmtid \\ }
\fi

\maketitle

% --- Abstract ---
\begin{abstract}
The abstract goes here.
\end{abstract}

\IEEEpeerreviewmaketitle

% --- Sections ---
\section{Introduction}
SIBGRAPI 2020 \cite{Sibgrapi2020}.

\subsection{Subsection Heading Here}
Subsection text here.

\subsubsection{Subsubsection Heading Here}
Subsubsection text here.

% --- Figure Example ---
\begin{figure}[!t]
\centering
\includegraphics[width=2.5in]{sibgrapi2020.png}
\caption{SIBGRAPI - Conference on Graphics, Patterns and Images.}
\label{fig_sim}
\end{figure}

% --- Double Figure Example ---
\begin{figure*}[!t]
\centering
\subfloat[Case I]{\includegraphics[width=2.5in]{sibgrapi2020.png}%
\label{fig_first_case}}
\hfil
\subfloat[Case II]{\includegraphics[width=2.5in]{sibgrapi2020.png}%
\label{fig_second_case}}
\caption{SIBGRAPI - Double figure example.}
\label{fig_sim2}
\end{figure*}

% --- Table Example ---
\begin{table}[!t]
\renewcommand{\arraystretch}{1.3}
\caption{An Example of a Table}
\label{table_example}
\centering
\begin{tabular}{|c||c|}
\hline
One & Two\\
\hline
Three & Four\\
\hline
\end{tabular}
\end{table}

\section{Conclusion}
The conclusion goes here.

% --- Acknowledgment ---
\section*{Acknowledgment}
The authors would like to thank...

% --- Bibliography ---
\bibliographystyle{IEEEtran}
\bibliography{example}

\end{document}